% !PRINTED 399
\weberpage{399}
\begin{center}
\S.\,128.

\gesperrt{Ableitung aller Lösungen der Pell'schen Gleichung
aus der kleinsten positiven.}
\end{center}

Ist $t$, $u$ irgend eine Lösung der Pell'schen Gleichung, so
ist der Ausdruck
\[
\frac{t^2 - D u^2}{4},
\]
der den Werth $\pm 1$ hat, die Norm der beiden conjugirten Zahlen
\[
\frac{t + u\sqrt{D}}{2}, \quad \frac{t - u\sqrt{D}}{2}.
\]
Wir wollen sie die zu $\sqrt{D}$ gehörigen \gesperrt{Einheiten} nennen.
Sind also $(t_1, u_1)$, $(t_2, u_2)$ irgend zwei (positive oder negative)
Lösungen der Pell'schen Gleichung
\[
t^2 - D u^2 = \pm 4,
\]
so haben die beiden Zahlen
\[
\Theta_1 = \frac{t_1 + u_1 \sqrt{D}}{2}, \quad \Theta_2 = \frac{t_2 + u_2 \sqrt{D}}{2}
\]
die Norm $\pm 1$, und das Gleiche gilt also auch von ihrem Product
\[
\Theta_3 = \Theta_1 \Theta_2 = \frac{t_3 + u_3 \sqrt{D}}{2}.
\]
Darin ist
\[
t_3 = \frac{t_1 t_2 + u_1 u_2 D}{2}, \quad u_3 = \frac{t_1 u_2 + t_2 u_1}{2},
\]
und $t_3$ und $u_3$ sind, wie man sofort sieht, ganze Zahlen.

Denn wenn $D$ gerade ist, so sind $t_1$, $t_2$ gerade, und wenn $D$
ungerade ist, so sind $t_1$, $u_1$ entweder beide gerade oder beide
ungerade, und ebenso $t_2$, $u_2$. Es ist also $\Theta_3$ gleichfalls eine zu
$\sqrt{D}$ gehörige Einheit, und es folgt, \gesperrt{dass das Product zweier}
(und folglich auch mehrerer) \gesperrt{Einheiten wieder eine Einheit ist.}

Ist $\Theta$ eine beliebige Einheit, so ist $\pm \Theta^{-1}$ die zu $\Theta$ conjugirte
Einheit. Es gehören unter den von $\pm 1$ verschiedenen
Einheiten immer vier zusammen,
\[
\pm \Theta, \quad \pm \Theta^{-1},
\]
unter denen \gesperrt{eine} positiv und grösser als 1 ist. Diese entspricht
%CONT>>
